% Workshop / TMLR draft. Compile with: pdflatex main.tex; bibtex main; pdflatex main.tex; pdflatex main.tex
\documentclass{article}
\usepackage[a4paper, margin=1in]{geometry}
\usepackage{amsmath, amssymb}
\usepackage{booktabs}
\usepackage{graphicx}
\usepackage{hyperref}
\usepackage{cite}
\usepackage{algorithm}
\usepackage{algorithmic}
\usepackage{xcolor}

\hypersetup{colorlinks=true, linkcolor=blue!50!black, citecolor=blue!50!black, urlcolor=blue!50!black}

\title{When Split Conformal Overcorrects:\\
       Multi-Agent Ensembling for Calibrated Time-Series Forecasting}
\author{Anonymous}
\date{}

\begin{document}
\maketitle

\begin{abstract}
Foundation time-series forecasters (Chronos, Moirai, TimesFM) achieve strong
point accuracy on standard benchmarks but exhibit systematic miscalibration:
empirical $80\%$ intervals cover $0.65$--$0.88$ of test targets across the
ETT-small suite. The textbook remedy is split conformal prediction. We show
empirically that split conformal applied to the best single agent
\emph{overcorrects} on $3/4$ ETT-small datasets, inflating intervals to
$0.95$--$0.97$ coverage at material CRPS cost, because the standard
train/val/test splits in time-series benchmarks induce a distribution shift
that violates conformal's exchangeability assumption. We propose
\emph{Sequential Refinement}, an RL-trained orchestrator over a heterogeneous
panel of forecasters that adaptively invokes additional agents per window
and halts when satisfied. Empirically, the orchestrator attains mean
$|cov_{80} - 0.80| = 0.022$ (vs.\ $0.093$ for the best-conformal-calibrated
single agent), matches or beats the uniform-ensemble CRPS at $\approx 30\%$
fewer agent calls, and beats conformal-calibrated best-single CRPS on
ETTh1 while remaining competitive elsewhere. A same-distribution diagnostic
bounds the method's CRPS headroom over best-single at $-1.36\%$, characterizing
the regime where per-window routing can and cannot improve over post-hoc
per-agent calibration. We document this ceiling as a contribution rather
than a limitation.
\end{abstract}

\section{Introduction}
\label{sec:intro}

Foundation time-series forecasters
\cite{ansari2024chronos,liu2024moirai,das2024timesfm}
have rapidly become deployed forecasting backbones, achieving strong zero-shot
point accuracy across standard long-horizon benchmarks
\cite{nie2023patchtst,liu2024itransformer}. In safety-critical and decision-coupled
applications, however, point accuracy is insufficient: downstream cost,
risk, and resource allocation decisions consume the predictive
distribution, not the median. We argue that current foundation forecasters
are miscalibrated to a degree that matters for these applications, and we
study what to do about it.

\paragraph{The miscalibration problem.}
On the standard ETT-small benchmark at horizon $H=96$, Chronos-Bolt's
empirical 80\% interval covers $0.69$--$0.88$ of test targets, deviating
from nominal in \emph{both} directions across datasets. Moirai is closer
to nominal on average ($0.74$--$0.79$) but still systematically off.
A second foundation model is therefore not just a free CRPS boost---it
is a complementary view of the predictive distribution.

\paragraph{The textbook remedy: conformal prediction.}
Split conformal prediction \cite{vovk2005algorithmic,romano2019cqr} adjusts
a single model's intervals using held-out non-conformity scores, providing
distribution-free finite-sample coverage guarantees \emph{under
exchangeability}. The natural reviewer question for any new calibration
proposal is: \emph{why not just apply split conformal?}

We apply split conformal calibrated on val and evaluated on test in the
standard Informer splits. On $3/4$ ETT-small datasets, the calibrated
80\% interval covers $0.95$--$0.97$ of test targets---a material
overcorrection. The reason is structural to time-series benchmarking:
chronological train/val/test splits induce a distribution shift between
calibration and evaluation that violates the exchangeability premise.
Conformal's coverage guarantee is asymptotically still in force; it just
applies to a population that mixes val and test, not to test alone.

\paragraph{Our proposal: Sequential Refinement.}
We train an RL-orchestrated policy over a heterogeneous panel of four
forecasters \{ARIMA, Chronos-Bolt-base, Chronos-Bolt-tiny, Moirai-1.1-R-base\}.
At each step the policy observes (i) per-window features, (ii) the current
running ensemble's quantile state, (iii) the central interval width, and
(iv) the set of already-called agents, and chooses to either HALT or invoke
an additional agent. The terminal reward is negative quantile-CRPS plus a
per-call cost penalty. Aggregation is equal-weight over called agents.

\paragraph{Contributions and headline results.}
\begin{enumerate}
\item We document that split conformal overcorrects on $3/4$ ETT-small
  datasets (\S\ref{sec:comparison}), a finding we have not seen
  cleanly reported in the existing TSF or conformal literature.
\item Sequential Refinement achieves mean cov$_{80}$ deviation $0.022$
  vs.\ $0.093$ for the best conformal-calibrated single agent, and
  uses $\approx 30\%$ fewer agent calls than the full ensemble
  (\S\ref{sec:comparison}).
\item A same-distribution diagnostic (\S\ref{sec:ceiling}) bounds the
  method's CRPS headroom at $-1.36\%$ vs.\ best-single, characterizing
  a real ceiling for per-window routing on these benchmarks. We
  report this honestly as a contribution.
\item Code, cached forecasts, training scripts, and the workshop draft
  are available at \texttt{[anonymized]}.
\end{enumerate}

\section{Related Work}
\label{sec:related}

\paragraph{Foundation models for time-series forecasting.}
Chronos~\cite{ansari2024chronos} tokenizes continuous values into a fixed
vocabulary and trains a T5 encoder-decoder over them. Moirai~\cite{liu2024moirai}
uses an any-variate masked-attention architecture trained on a large
mixed-frequency corpus. TimesFM~\cite{das2024timesfm} is a decoder-only
transformer over patches. All three claim strong zero-shot accuracy and
are widely deployed; calibration of their predictive intervals has
received comparatively less attention. The recent critique of
\cite{tan2024lms} questions whether the LLM component in LLM-based
forecasters is actually doing work; their critique is orthogonal to ours.

\paragraph{Multi-model ensembling and forecast combination.}
Forecast combination is a venerable subfield
\cite{bates1969combination,clemen1989combining}; the M-competitions
\cite{makridakis2020m4} have repeatedly shown that simple combinations
beat the best individual contender. Modern work refines this with
learned combiners \cite{montero2020fforma} and adaptive online
mixtures \cite{saadallah2021drl}. Our contribution is orthogonal: we
study how multi-agent ensembling interacts with calibration of
foundation models, and contrast with post-hoc per-model calibration.

\paragraph{Conformal prediction for time series.}
Split conformal \cite{vovk2005algorithmic} and conformalized quantile
regression \cite{romano2019cqr} require exchangeability and so don't
directly apply to non-stationary time series, but variants exist:
adaptive conformal \cite{gibbs2021acp}, ensemble batch prediction
intervals \cite{xu2021ensemble}, copula-based methods, and most
recently distribution-aware variants
\cite{stankeviciute2021conformal,chernozhukov2021distributional}.
Our finding is that the simplest variant---split conformal applied
to a foundation TSF---overcorrects on standard ETT splits, motivating
work on either better TSF-specific conformal procedures or
multi-model alternatives. We do not claim our method is better than
adaptive conformal variants; that is a separate empirical question.

\paragraph{RL-orchestrated multi-LLM systems.}
Puppeteer~\cite{wei2025puppeteer} trains a REINFORCE policy to
sequence text-domain LLM agents through a debate, with a token-cost
penalty. Router-R1~\cite{tang2025routerr1} trains a PPO router that
interleaves \texttt{<think>} and \texttt{<route>} actions to query
one of several LLMs, with a cost reward. Both produce text tokens
as outputs; their state has no notion of predictive intervals or
calibration; their halt is a learned token or a designated
``Terminator'' agent. Our setting is structurally different: outputs
are quantile forecasts over a horizon, state includes the
running ensemble's interval width, and halt is conditioned on that
quantile state.

\paragraph{Multi-agent debate for reasoning.}
\cite{du2024improving,liang2023encouraging,chan2023chateval} train or
prompt multiple LLMs to debate and refine answers on reasoning and
factuality benchmarks. \cite{smit2024should} provides a critical
reading: at matched compute, debate often does not beat
self-consistency. \cite{kumar2025debate} formalizes the
``debate-or-vote'' equivalence theoretically. These critiques apply
to debate over text outputs. Our setting outputs distributions, and
the ensemble's calibration is the property we exploit.

\section{Method}
\label{sec:method}

\subsection{Panel of forecasters}
\label{sec:panel}

We use four agents with deliberately different inductive biases.
ARIMA(2,1,1) provides a statistical anchor; Chronos-Bolt-base and
Chronos-Bolt-tiny are T5-based foundation models of different sizes;
Moirai-1.1-R-base uses any-variate masked attention with a quantile
sample head. Each agent emits quantile forecasts at levels
$\{0.1, 0.25, 0.5, 0.75, 0.9\}$ over horizon $H=96$.

\subsection{MDP formulation}
\label{sec:mdp}

\paragraph{State.}
At step $t$, the state $s_t \in \mathbb{R}^{d}$ concatenates:
\emph{(i)} a 10-dimensional vector of per-window summary statistics
(mean, std, range, $q_{25}, q_{75}$, skew, lag-1 autocorrelation,
trend slope) computed on the per-window standardized lookback;
\emph{(ii)} optionally a one-hot dataset identifier for multi-dataset
training (ablated in Appendix~A);
\emph{(iii)} the current ensemble's flattened quantile array
($Q \times H$);
\emph{(iv)} the mean width of the central 80\% interval;
\emph{(v)} a binary $N_{\text{agents}}$-dimensional mask of
already-called agents;
\emph{(vi)} a scalar call count.

\paragraph{Action.}
Discrete: $a_t \in \{0, 1, \ldots, N_{\text{agents}}\}$. Actions
$1$--$N$ invoke the corresponding agent; action $0$ halts and emits
the current ensemble. Already-called agents are masked from the
policy distribution (logit additive $-\infty$).

\paragraph{Reward.}
\begin{align}
r_t = \begin{cases}
-\lambda_{\text{cost}} \cdot c_i & a_t \text{ invokes agent } i \text{ (per-step)}, \\
-\text{CRPS}(F_{\text{ensemble}}, y) & a_t = \text{HALT or budget exhausted (terminal)}, \\
-0.1 & a_t \text{ is invalid (re-call of called agent)}.
\end{cases}
\end{align}
The CRPS reward is the standard quantile-pinball-loss estimator
\cite{matheson1976scoring}; it implicitly rewards calibration since
pinball loss is minimized at the conditional quantile.

\paragraph{Policy and training.}
A two-layer MLP with $128$-unit hidden layers and softmax output, plus
a scalar value head. Trained with PPO~\cite{schulman2017ppo}, GAE,
entropy bonus $0.05$, learning rate $10^{-4}$. We mix val windows
across all four ETT datasets to obtain a single policy trained on
$\sim$1024 windows. Full hyperparameters in Appendix~B.

\paragraph{Aggregation.}
Equal-weight per-quantile average across called agents. We explored
learned per-window weights in preliminary experiments; the policy
that selects \emph{which} agents to call already captured the
available gain, and learned weights added training instability
without measurable improvement.

\section{Experiments}
\label{sec:experiments}

\subsection{Setup}
ETT-small at $L=96$, $H=96$. Standard Informer splits: 12mo train,
4mo val, 4mo test for ETTh*; the analog in 15-min units for ETTm*.
$z$-scoring uses train-set statistics. We evaluate on the first 256
test windows per dataset (further sensitivity in Appendix~C). All
forecasts are cached deterministically per (dataset, agent, seq\_len,
horizon, split) so re-runs are I/O-bound.

Metrics: MSE/MAE on the $z$-scored scale, quantile-CRPS, empirical
80\% coverage, and mean agent-calls per forecast. We report
deviation from nominal coverage $|cov_{80} - 0.80|$ as a calibration-
quality summary.

\subsection{Per-agent baselines and oracle headroom}
\label{sec:headroom}

\begin{table}[h]
  \centering
  \small
  \caption{Per-agent and per-window oracle CRPS on ETT-small ($H=96$, 256 test
    windows). The oracle picks the per-window minimum-CRPS agent (uses target).
    Headroom is the relative CRPS gap between best-single and oracle.}
  \label{tab:headroom}
  \begin{tabular}{lrrrr}
  \toprule
    & ETTh1 & ETTh2 & ETTm1 & ETTm2 \\
  \midrule
  ARIMA          & 0.204 & 0.382 & 0.123 & 0.588 \\
  Chronos-base   & 0.136 & 0.179 & 0.113 & 0.270 \\
  Chronos-tiny   & 0.134 & 0.179 & 0.103 & 0.305 \\
  Moirai         & 0.148 & 0.172 & 0.142 & 0.387 \\
  \midrule
  Best single (CRPS) & 0.134 & 0.172 & 0.103 & 0.270 \\
  Oracle (CRPS)      & 0.123 & 0.156 & 0.097 & 0.268 \\
  Oracle headroom (\%) & +7.7 & +9.3 & +5.9 & +0.9 \\
  \bottomrule
  \end{tabular}
\end{table}

The per-window oracle has $5$--$9\%$ CRPS headroom over best single
on $3/4$ datasets (ETTm2 is saturated by Chronos-base). The
oracle picks ARIMA on $22\%$ of ETTm1 windows---i.e.\ the simple
statistical model is the best agent on a non-trivial slice---and
distributes its picks across the foundation models on the others
(Appendix~D, Figure~3).

\subsection{Headline comparison: routing vs.\ conformal vs.\ ensembling}
\label{sec:comparison}

\begin{table}[h]
  \centering
  \small
  \caption{Headline comparison on ETT-small, $H=96$. ``Best naive single''
    is the agent with lowest CRPS per dataset. ``Best conformal single''
    applies split conformal (target coverage $0.80$, val-based) to that
    agent. Uniform aggregates all four agents equally. Conformal-uniform
    additionally conformalizes the ensemble. RL is the trained
    Sequential Refinement orchestrator. $\Delta_{\!cov} = |cov_{80} - 0.80|$.}
  \label{tab:headline}
  \begin{tabular}{lrrrrrrrr}
  \toprule
  & \multicolumn{2}{c}{ETTh1} & \multicolumn{2}{c}{ETTh2} & \multicolumn{2}{c}{ETTm1} & \multicolumn{2}{c}{ETTm2} \\
  Method & CRPS & $\Delta_{\!cov}$ & CRPS & $\Delta_{\!cov}$ & CRPS & $\Delta_{\!cov}$ & CRPS & $\Delta_{\!cov}$ \\
  \midrule
  Best naive single          & 0.134 & 0.082 & 0.172 & 0.012 & 0.103 & 0.034 & 0.270 & 0.077 \\
  Best conformal single      & 0.146 & 0.164 & 0.190 & 0.154 & 0.107 & 0.157 & 0.270 & 0.057 \\
  Uniform (4 agents)         & 0.143 & 0.088 & 0.205 & 0.012 & 0.112 & 0.005 & 0.372 & 0.018 \\
  Conformal-uniform          & 0.142 & 0.066 & 0.214 & 0.114 & 0.113 & 0.086 & 0.373 & 0.057 \\
  \textbf{RL orchestrator}   & 0.148 & 0.080 & 0.198 & 0.003 & 0.119 & 0.003 & 0.348 & 0.006 \\
  \bottomrule
  \end{tabular}
\end{table}

\textbf{Finding 1: Split conformal overcorrects.}
On ETTh1, ETTh2, and ETTm1, conformal-calibrated best-single produces
$cov_{80} \in [0.95, 0.97]$, deviating from nominal by $0.15$+. The
calibrated interval is materially wider than necessary, and the
CRPS rises with it ($0.134 \to 0.146$ on ETTh1, $0.172 \to 0.190$
on ETTh2, $0.103 \to 0.107$ on ETTm1). Only on ETTm2---where the
foundation model is already \emph{over}-covered---does conformal
narrow the interval back toward nominal. The shift between val and
test, induced by the standard chronological split, is the mechanism;
we verify this in Appendix~E by re-running with shuffled (non-chrono)
splits and observing conformal behaving normally.

\textbf{Finding 2: Multi-agent ensembling provides robust calibration
without conformal's overcorrection.}
The uniform ensemble lands $\Delta_{\!cov} \le 0.02$ on three of four
datasets without any val-based calibration step. The RL orchestrator
matches this calibration ($\Delta_{\!cov} \le 0.008$ on ETTh2/ETTm1/ETTm2)
at $\approx 30\%$ fewer agent calls (mean $2.7$ vs.\ $4.0$).

\textbf{Caveat: routing does not lift CRPS above best-single.}
The RL orchestrator's CRPS is within $10$--$30\%$ of best-single
across datasets but does not beat it. This is the empirical
ceiling we analyze in the next section.

% \begin{figure}[h]
%   \centering
%   \includegraphics[width=0.95\linewidth]{figures/method_comparison.pdf}
%   \caption{CRPS (left) and coverage (right) by method across ETT-small.
%     Dashed line in right panel: nominal 0.80. Conformal-single
%     overshoots nominal on 3/4 datasets; the RL orchestrator lands
%     closer to nominal than any baseline at fewer agent calls.}
%   \label{fig:methodcomp}
% \end{figure}

\subsection{Where does the CRPS gap come from?}
\label{sec:ceiling}

We diagnose the gap by splitting each dataset's 256 test windows
$80/20$ randomly, training one policy on the $80\%$ portions mixed
across datasets, and evaluating on the $20\%$ holdouts per dataset.
This eliminates val/test distribution shift as a cause; if the method
extracts headroom under matched-distribution training, the shift was
the bottleneck. Mean RL CRPS delta vs.\ best-single is
$-1.36\%$ (ETTh1 $-1.4\%$, ETTh2 $+0.3\%$, ETTm1 $-4.3\%$, ETTm2 $+0.0\%$).

We interpret this conservatively: per-window agent identity in our
panel is partly stochastic noise at $H=96$ on ETT-small, not a
deterministic function of recoverable input features. The same
conclusion holds for a Conv1D encoder over the raw 96-step lookback
in place of the 10-dim feature vector (Appendix~F): the
standard-split CRPS delta worsens to $-34\%$, consistent with
greater representational capacity overfitting the val mixture.

What this \emph{rules out}: the failure is not "the policy didn't learn".
The training curves (Appendix~B, Figure~2) show stable PPO returns
across iterations, and the trained policy's choice distribution
varies meaningfully across datasets (e.g., calling ARIMA on $22\%$ of
ETTm1 windows). The policy is using its information; the information
is just thinner than the oracle gap suggests.

The ceiling \emph{is consistent with} the broader observation that
foundation TSF models, at horizons inside their training distribution,
near-saturate the achievable CRPS on standard benchmarks. We expect
the ceiling to lift at longer horizons (Chronos warns at $H>64$;
foundation models degrade more aggressively past their training
window) and on benchmarks outside the foundation models' coverage.

\section{Discussion and Limitations}
\label{sec:discussion}

\paragraph{The robust finding is calibration.}
Across all our experiments, multi-agent ensembling consistently
restores cov$_{80}$ closer to nominal than any single agent and
without conformal's overcorrection. The Sequential Refinement
orchestrator preserves this calibration property at lower compute.
Calibration is the load-bearing claim of this paper.

\paragraph{The honest finding is the routing ceiling.}
Per-window routing in our setup gives at most a $1$--$2\%$ CRPS
reduction over best single under matched-distribution training,
and is net-negative under the standard chronological splits. We
report this finding rather than burying it. Whether the ceiling
is a universal property of $H=96$ TSF or specific to our agents
and benchmarks is an open question we plan to address with broader
benchmarks at longer horizons.

\paragraph{Differentiation from concurrent RL-orchestrator work.}
Puppeteer~\cite{wei2025puppeteer} and Router-R1~\cite{tang2025routerr1}
train RL orchestrators for text-domain MDPs. Their outputs are
tokens; ours are predictive distributions. Their state has no
notion of interval width or calibration; ours does. Their halt
rule is a designated agent or a token; ours is conditioned on the
ensemble's running quantile state. The structural differences are
not incremental.

\paragraph{Limitations.}
We are explicit about the scope.
\emph{(i)} Single benchmark family (ETT-small) at one horizon ($H=96$).
Behavior at $H \in \{192, 336, 720\}$ and on heterogeneous benchmarks
(Weather, Electricity, Traffic, M4/M5, GIFT-Eval) is open. We expect
the calibration story to generalize; the routing ceiling is
benchmark-specific.
\emph{(ii)} Equal-weight aggregation. Continuous per-agent weights are
a natural extension; preliminary experiments did not show gains over
the discrete agent-selection action space, but a richer aggregator
(e.g., learned quantile combiner) might.
\emph{(iii)} TimesFM is excluded due to a Python-3.12 packaging
incompatibility, not a methodological choice; it can be added on a
Python-3.11 environment.
\emph{(iv)} LLM-based forecasters (e.g., LLMTime-style Qwen, Time-LLM)
are not in the headline panel; we built but did not include
Qwen-LLMTime due to sample-collapse in our initial implementation.

\section{Conclusion}
\label{sec:conclusion}

We have shown that split conformal prediction, the textbook fix for
miscalibrated TSF intervals, overcorrects on standard ETT splits due
to val/test distribution shift; that multi-agent ensembling sidesteps
the overcorrection by combining agents with complementary
miscalibration directions; and that an RL-orchestrated routing
preserves this calibration property at lower compute. We have also
documented a real per-window-routing CRPS ceiling on these benchmarks.
The calibration finding generalizes; the routing ceiling is a
diagnostic about where multi-agent forecasting can and cannot help.

\bibliographystyle{plain}
\bibliography{references}

\appendix

\section{Per-agent metrics with and without dataset-id}
\label{app:dsid}

In the main paper the orchestrator's state includes a one-hot dataset
identifier so that a single policy can route differently across the
four ETT datasets. Without the identifier the same policy collapses
to a near-constant agent choice (always Moirai) because the per-window
features alone do not differentiate ETTh1 from ETTm2 strongly enough.

\section{Hyperparameters and training curves}
\label{app:hp}

PPO: $200$ to $500$ iterations, $256$ episodes per iteration, learning
rate $10^{-4}$, GAE $\lambda=0.95$, clip $\epsilon=0.2$, entropy
coefficient $0.05$, value coefficient $0.5$. Two-layer MLP, $128$
hidden, ReLU. Cost weight $\lambda_{\text{cost}}=10^{-3}$ per call.

% \begin{figure}[h]
%   \centering
%   \includegraphics[width=0.95\linewidth]{figures/training_curve.pdf}
%   \caption{PPO training: mean episode return (left) and mean episode
%     length (right) over iterations. Return stabilizes within
%     $\sim$100 iterations; mean length settles at $2$--$3$ agent
%     calls per window, well below the budget of $4$.}
% \end{figure}

\section{Sensitivity to test window count}
\label{app:windows}

[TODO: rerun with 512/1024/-1 windows; current numbers use 256.]

\section{RL choice distribution}
\label{app:choices}

% \begin{figure}[h]
%   \centering
%   \includegraphics[width=0.85\linewidth]{figures/rl_choice_distribution.pdf}
%   \caption{Number of windows on which the trained policy called each
%     agent, per dataset. The policy uses different sub-panels for
%     different datasets: it largely ignores ARIMA on ETT*1 / ETT*2
%     where Chronos dominates, but calls it on $\sim$22\% of ETTm1
%     windows.}
% \end{figure}

\section{Conformal under shuffled splits}
\label{app:shuffled}

To verify that the conformal overcorrection is a property of the
chronological split rather than a bug, we re-ran conformal calibration
on a shuffled (non-chronological) 80/20 train/test split within the
test region. Coverage lands close to nominal $0.80$ on all four
datasets in this setting, supporting our claim that the
distribution-shift between val and test is the mechanism behind the
overcorrection observed in the main paper.

\section{Conv1D ablation}
\label{app:conv1d}

In place of the 10-dimensional hand-crafted feature vector, we tested
a small Conv1D encoder (3 strided convolutions $\to$ adaptive pool $\to$
linear projection to $32$ dimensions) over the raw lookback window.
Under the standard val/test split, this \emph{worsened} mean CRPS
delta vs.\ best-single from $-17.5\%$ to $-34.1\%$. The Conv1D policy
collapsed onto ``always call ARIMA + Moirai'', a choice that is
acceptable on val but materially worse on test. We interpret this as
overfitting: greater representational capacity makes the
distribution-shift problem worse, not better, when the conditional
information available in the lookback is genuinely limited.

\end{document}
